\chapter{Beyond the Blue Book}
\label{ch:beyond}

\standfirst{Everything this language has that Smalltalk-80 did not, what it
deliberately does not have, and how to turn it on. If you already know
Smalltalk, this is the chapter to read first.}

\section{The dialect switch}

The compiler has two dialects. \ct{bluebook} is the 1983 grammar and the 1983
block semantics; \ct{closures} is everything in this chapter. Every profile in
\ct{profiles/} except \ct{bluebook} itself selects \ct{closures}, with one line:

\begin{st}
#dialect : 'closures',
\end{st}

Building with \ct{-manifest sources/MANIFEST} and no profile gives you the
museum piece: no closures, no exceptions, no \ct{Mutex}.

\begin{stnote}[A doit is compiled in the profile's dialect]
\ct|#dialect| governs the packages a profile loads, and it governs the
expression you hand to \ct{-eval}, \ct{-startup} or \ct{-serve} and each
example \ct{-doctests} runs. So \ct{[:x | x] class} answers
\ct{BlockClosure} under \ct{profiles/st2026.profile} and \ct{BlockContext}
under \ct{profiles/bluebook.profile}, which is what each profile says it is.

There is only one code generator, so a method you accept in the Browser and
an expression you evaluate in a workspace are compiled by the same one that
built the image. The image reaches it through primitive 228; the 1983 parser
is what pretty-prints, decompiles and answers \ct{parseSelector:}.
\end{stnote}

\section{Full closures}

The largest of the differences, and the one everything else rests on. Blocks
capture their enclosing variables, survive the method that made them, and may
be entered more than once at the same time---which matters here, because two
workers may evaluate the same block object concurrently.

\begin{st}
makeCounter
    | n |
    n := 0.
    ^[n := n + 1]           "still works after makeCounter has returned"
\end{st}

Non-local return (\ct|^| inside a block), \ct{ensure:} and
\ct{ifCurtailed:} come with them, and so does the exception system built on
top---see Chapters~\ref{ch:blocks} and~\ref{ch:exceptions}.

The bytecodes are Squeak's V3PlusClosures set, not Pharo's SistaV1. That is a
free choice here because this system compiles Pharo's source itself rather
than loading Pharo's compiled methods; the only code that could tell the
difference is code that inspects bytecodes.

\section{Dynamic arrays}

\begin{st}
{ 1. 2 + 3. Date today }
\end{st}

Braces, expressions separated by periods, evaluated at run time. Compiled as
\ct{Array new:} filled by \ct{at:put:}, so it needed no new bytecode.

The contrast with \ct|#(1 2 3)| is the point: a literal array holds
compile-time constants and is one shared object; a dynamic array is built
fresh each time from whatever the expressions answer. It is also the only way
to get \ct{true}, \ct{false} and \ct{nil} into an array in this
system---inside \ct|#(...)| they are bare words and bare words become
symbols.

\section{Byte arrays}

\begin{st}
#[1 2 255]
#(#[1 2] #[3 4])        "nested inside a literal array; also fine"
\end{st}

Elements 0 to 255, a real \ct{ByteArray}.

\section{Block-local temporaries}

\begin{st}
[:x | | doubled | doubled := x * 2. doubled + 1] value: 3      "7"
[ | t | t := 10. t squared ] value                             "100"
\end{st}

Each temporary gets its own frame slot and is set to \ct{nil} at every
activation, so a block called twice does not see the last call's values.

\section{Pragmas}

\ct{<primitive: 60>} was the Blue Book's one annotation. It is now general:
any keyword message with literal arguments, several per method.

\begin{st}
annotated
    <test>
    <deprecated: 'use #normal'>
    <shared: #serialize>
    ^#annotated
\end{st}

Pragmas are attached to the compiled method and readable back:

\begin{st}
(Unwind class >> #annotated) pragmas size                        "3"
(Unwind class >> #annotated) hasPragma: #shared:                 "true"
((Unwind class >> #annotated) pragmaAt: #deprecated:) arguments first
                                                     "'use #normal'"
\end{st}

The two primitive forms still mean something to the compiler; every other
pragma is stored and otherwise ignored, which is exactly what you want---the
meaning belongs to whatever framework reads them.

\begin{stnote}[One real limit]
Binary selectors are scanned greedily up to two characters, so a closing
\ct{>} immediately followed by another binary character arrives as one token
and the pragma is not recognised. In practice pragmas end lines and this never
comes up, but the limit is real.
\end{stnote}

\subsection{Named primitives}

\begin{st}
<primitive: 'primitiveName' module: 'ModuleName'>
\end{st}

Parsed and compiled as Squeak's primitive 117 with the descriptor as literal
zero. The VM does not yet dispatch it, so the Smalltalk body always runs. It
is here so that Pharo source containing them can be loaded.

\section{Traits}

A trait is a named bag of methods with no instances and no place in the
superclass chain. This system applies them by \emph{flattening at load
time}---the trait's source is compiled into each class that uses it, which is
why an instance-variable reference inside a trait method means the same thing
in every class that takes it.

\begin{st}
Trait {
	#name : 'TGreeting',
	#category : 'Probe-Core'
}

{ #category : 'greeting' }
TGreeting >> greeting [
	^'hello from ', self subject
]
\end{st}

\begin{st}
Class {
	#name : 'Greeter',
	#superclass : 'Object',
	#traits : 'TGreeting',
	#classTraits : 'TGreeting classTrait',
	#category : 'Probe-Core'
}

{ #category : 'greeting' }
Greeter >> subject [
	^'a class'
]
\end{st}

\ct{Greeter new greeting} answers \ct{'hello from a class'}: the trait supplied
\ct{greeting}, the class supplied \ct{subject}, and the class's own method
always wins over the trait's.

Composition supports \ct{+} only. Trait subtraction (\ct{-}) and aliasing
(\ct{@}) are rejected with a report rather than silently mis-loaded, and there
is no trait reflection or update propagation: flattening is a load-time
operation, and changing a trait afterwards does not change the classes that
took it.

\section{\texttt{new} calling \texttt{initialize}}

Pharo classes assume \ct{new} sends \ct{initialize}; the 1983
\ct{Behavior>>new} is a primitive and nothing else. Rather than change
\ct{Behavior>>new}---which would double-initialise the thirty-four 1983
classes that already write \ct|^super new initialize| by hand---the loader
synthesises the 1983 idiom for any loaded class that has \ct{initialize}
reachable and no class-side \ct{new} of its own, and reports every class it
did that to.

Classes you type into a Browser get no such help. Write the constructor.

\section{What is deliberately missing}

\subsection{Scaled decimals}

\ct{1.23s2} is not implemented and will not be, and the reason is worth
knowing. The other five extensions were clean \emph{syntax errors} before, so
adding them could not take meaning away from anything. \ct{1.23s2} already
parses: a unary selector needs no space before it, so \ct{3factorial} is
\ct{3 factorial} and \ct{1.23s2} is \ct{1.23 s2} by the same rule. It is the
only one of the six with an existing meaning to take away, there is no
\ct{ScaledDecimal} class for it to answer, and \ct{Float} does not implement
\ct{s2}, so you get a doesNotUnderstand and \ct{nil}.

\subsection{The tool ecosystem}

No Morphic, Bloc, Spec or Athens; no Iceberg, Metacello or Monticello; no
Fuel; no Opal or RB-AST. The first group is single-threaded by construction.
The second is replaced by directories of Tonel files and a profile. The last
assumes SistaV1 bytecodes and Spur literals.

\section{Summary, for people arriving from Pharo}

\begin{center}
\small
\begin{tabular}{@{}lp{0.52\textwidth}@{}}
\toprule
closures, \ctp{ensure:}, exceptions & yes, in the \ctp{closures} dialect\\
\ct|{ }|, \ct|#[ ]|, block temps, pragmas & yes\\
traits & yes, flattened at load, \ctp{+} only\\
\ctp{new} sends \ctp{initialize} & for loaded packages, synthesised
  and reported; not automatic for classes you type in\\
\ctp{1.23s2} & no, and it silently answers \ctp{nil}\\
\ct|#(true)| & a Symbol, not \ctp{true}\\
\ct|#foo printString| & \ct|'foo'|, not \ct|'#foo'|\\
\ctp{Ctrl-D}, \ctp{Ctrl-P} & no; \ctp{do it} and \ctp{print it} are in the
  middle-button menu\\
same-priority atomicity & \textbf{gone}; see Chapter~\ref{ch:contract}\\
\bottomrule
\end{tabular}
\end{center}
